% Clase 22 — un circuito magnético: núcleo cerrado de alta permeabilidad con un
% bobinado de N vueltas, tramos de materiales distintos y un entrehierro chico.
% Las líneas de B circulan por el núcleo; la curva C va por el medio del núcleo.
\begin{tikzpicture}[scale=1]

  % nucleo (anillo rectangular)
  \fill[figgray!30, even odd rule]
    (0,0) rectangle (6,4)  (1,1) rectangle (5,3);
  \draw[figline] (0,0) rectangle (6,4);
  \draw[figline] (1,1) rectangle (5,3);
  % tramo de otro material (B) arriba
  \fill[figamber!30] (2.4,3) rectangle (3.6,4);
  \draw[figamber, line width=0.8pt] (2.4,3) -- (2.4,4) (3.6,3) -- (3.6,4);
  \node[figlbl, fill=figamber!30, inner sep=1pt] at (3.0,3.72) {B};
  \node[figlbl] at (5.78,2.85) {A};
  % entrehierro en el lado inferior
  \fill[white] (2.85,-0.05) rectangle (3.15,1.05);
  \draw[figline] (2.85,0) -- (2.85,1) (3.15,0) -- (3.15,1);
  \node[figlblaux, anchor=north] at (3.0,-0.1) {entrehierro};

  % bobinado en el lado izquierdo
  \foreach \yy in {1.2,1.5,1.8,2.1,2.4,2.7} {
    \draw[figred, line width=1pt] (-0.12,\yy) -- (1.12,\yy+0.12);
  }
  \draw[figred, line width=1pt] (-0.12,1.2) -- (-0.7,1.2);
  \draw[figred, line width=1pt] (-0.12,2.82) -- (-0.7,2.82);
  \node[figlbl, text=figred, anchor=east] at (-0.75,2.0) {$N,\ I$};

  % curva C por el medio
  \draw[figblue, dashed, line width=0.8pt] (0.5,0.5) rectangle (5.5,3.5);
  \draw[figvec] (5.5,1.6) -- (5.5,2.3);
  \node[figlbl, text=figblue] at (4.3,0.25) {$C$};

\end{tikzpicture}
